

%
%% ===========================================================
%\usepackage[tikz]{bclogo}
%
%\renewcommand\logowidth{20pt}
%
%%\newcommand\logoPython{\includegraphics[width=17pt]{include/logoPython.pdf}}
%\newenvironment{python}{%
%	%	\tikzset{external/export=false}%
%	\begin{center}%
%		\begin{minipage}{1\textwidth}%
%			\begin{bclogo}[logo=\bccrayon, arrondi=0.1, couleur=gray!20, barre = snake, epBarre = 1.0, ombre=true, epOmbre = 0.2, couleurOmbre = black, tailleOndu=1.5]{Et en Python?}%
%				%
%			}{\end{bclogo}%
%		\end{minipage}%
%	\end{center}%
%	%	\tikzset{external/export=true}%
%}%
%\newcommand{\outsidefootnote}[1]{\tcblower \footnotemark #1}
%% ===========================================================


\lstset{
	language=Python,
	frame=leftline,
	framerule=0.2mm,
	basicstyle=\footnotesize\ttfamily,
	breaklines=true,
	postbreak=\raisebox{0ex}[0ex][0ex]{\ensuremath{\color{red}\hookrightarrow\space}},
	%
	backgroundcolor = \color{lightgray!10},
	keywordstyle=\bfseries,
	commentstyle=\color{teal}\itshape,
	%
	numbers=left,
	numberstyle=\tiny,
	numbersep=5pt,
	xleftmargin = 1em,
	% framexleftmargin = 1em,
	%
	literate=% pour avoir les accents
	{é}{{\'e}}{1}%
	{è}{{\`e}}{1}%
	{à}{{\`a}}{1}%
	{ç}{{\c{c}}}{1}%
	{œ}{{\oe}}{1}%
	{ù}{{\`u}}{1}%
	{É}{{\'E}}{1}%
	{È}{{\`E}}{1}%
	{À}{{\`A}}{1}%
	{Ç}{{\c{C}}}{1}%
	{Œ}{{\OE}}{1}%
	{Ê}{{\^E}}{1}%
	{ê}{{\^e}}{1}%
	{î}{{\^i}}{1}%
	{ô}{{\^o}}{1}%
	{û}{{\^u}}{1}%
	{ë}{{\¨{e}}}1
	{û}{{\^{u}}}1
	{â}{{\^{a}}}1
	{Â}{{\^{A}}}1
	{Î}{{\^{I}}}1
}
